% !TEX root = ../notes.tex

\documentclass[../notes.tex]{subfiles}

\begin{document}

\section{April 6}
Here we go.

\subsection{\'Etale Cohomology}
\'Etale cohomology arises by changing the notion of ``open subset'' in algebraic geometry from (Zariski) open embeddings to \'etale morphisms. Thus, we should start by defining what it means for a morphism to be \'etale. Roughly speaking, an \'etale morphism should be a local diffeomorphism.
\begin{defihelper}[\'etale] \nirindex{etale@\'etale}
	Fix a morphism $f\colon Y\to X$ of schemes, and choose a point $y\in Y$.
	\begin{listalph}
		\item If $f$ takes the form $\Spec A[t_1,\ldots,t_n]/(g_1,\ldots,g_n)\to\Spec A$, then $f$ is \textit{\'etale} at $y$ if and only if the Jacobian matrix $\left[\del g_i/\del t_j\right]_{ij}$ is an invertible matrix with entries in $\OO_{Y,y}/\mf m_y$.
		\item In general, we say that $f$ is \textit{\'etale} at $y$ if and only if there are open neighborhoods $Y'\subseteq Y$ of $y$ and $X'\subseteq X$ of $f(y)$ such that the morphism $Y'\to X'$ is isomorphic to the situation in (a). Explicitly, $X'$ should be some affine scheme $\Spec A$, and $Y'$ should take the form $\Spec A[t_1,\ldots,t_n]/(g_1,\ldots,g_n)$, and the Jacobian at $y$ should be invertible.
	\end{listalph}
	We say that $f$ is \textit{\'etale} if and only if it is \'etale at every point $y\in Y$.
\end{defihelper}
\begin{remark}
	The implicit function theorem tells us that (a) is equivalent to $f$ being 
\end{remark}
\begin{remark}
	It is not totally obvious that (a) and (b) are equivalent. Namely, it is not obvious that one cannot pass to some suitable open subsets of $\Spec A$ and $\Spec A[t_1,\ldots,t_n]/(g_1,\ldots,g_n)$ to force the Jacobian to be invertible when it wasn't before. This doesn't happen, and the best argument gives a more intrinsic definition of being \'etale (e.g., being finitely presented, flat, and unramified).
\end{remark}
Now that we know what our new open subsets are, we would like to build a sheaf theory. This requires a notion of an open covering, for which our definition is unsurprising.
\begin{defihelper}[\'etale covering]
	An \textit{\'etale covering} of a scheme $X$ is a collection of \'etale morphisms $\{f_i\}$ with target $X$ for which $\bigcup_i\im f_i=X$ on the underlying topological spaces.
\end{defihelper}
\begin{definition}[small \'etale site]
	Fix a scheme $X$. Then the \textit{small \'etale site} of $X$ is the category $\op{\acute Et}(X)$ whose objects are \'etale morphisms $U\to X$ and morphisms are morphisms over $X$. We equip this category with coverings given by the collection of \'etale coverings: namely, a covering of an \'etale map $U\to X$ is just an \'etale covering $\{V_i\to U\}$ of $U$ by \'etale morphisms $V_i\to X$.
\end{definition}
\begin{remark}
	Part of the coherence here is that if we have two morphisms $X\to Y\to Z$ where $X\to Z$ and $Y\to Z$ are \'etale, then $X\to Y$ is \'etale.
\end{remark}
With a notion of a covering, we can produce a sheaf theory.
\begin{defihelper}[\'etale sheaf] \nirindex{etale sheaf@\'etale sheaf}
	Fix a scheme $X$. Then an \textit{\'etale presheaf} $\mf F$ on $X$ is a presheaf on $\op{\acute Et}(X)$. An \textit{\'etale sheaf} $\mc F$ on $X$ is a presheaf satisfying the sheaf condition; explicitly, for any open covering $\{V_i\}$ of some \'etale open $U\to X$, there is an equalizer diagram
	\[\mc F(U)\to\prod_i\mc F(V_i)\rightrightarrows\prod_{i,j}\mc F(V_i\times_XV_j).\]
	We let $\op{Sh}(X_{\mathrm{\acute et}})$ denote the category of abelian sheaves on $X_{\mathrm{\acute et}}$.
\end{defihelper}
\begin{example}
	There is an \'etale sheaf $\mathbb G_m$ on $X$ which sends $V\to X$ to $\OO(V)^\times$. We will not check that this satisfies the sheaf condition.
\end{example}
\begin{remark}
	It turns out that $\op{Sh}(X_{\mathrm{\acute et}})$ is an abelian category. The proof of this is similar to the Zariski case, and it proceeds by defining a notion of sheafification.
\end{remark}
With a sheaf theory, we also have a cohomology theory.
\begin{defihelper}[\'etale cohomology]
	Fix a scheme $X$. Then the functor $\Gamma\colon\op{Sh}(X_{\mathrm{\acute et}})\to\mathrm{Ab}$ given by taking global sections is left-exact, so we define its right-derived functors by $\mathrm H^i_{\mathrm{\acute et}}(X;-)$. These are the \textit{\'etale cohomology} functors.
\end{defihelper}
We close our discussion of this example.
\begin{example}
	Let $X$ be a field $\Spec k$. Then an \'etale morphism $Y\to X$ must be a zero-dimensional scheme, so the map $Y\to X$ can be seen to force $Y$ to be affine. In fact, using the local criterion, we conclude that $Y$ must be a disjoint union of finite separable extensions of $k$. Now, for an \'etale sheaf $\mc F$ on $X$, we define
	\[\mc F(k^{\mathrm{sep}})\coloneqq\colim_{L/k}\mc F(L),\]
	where the colimit is taken over finite separable extensions $L$ of $k$. For each Galois extension $L$ of $k$, we see that $\mc F(L)$ admits an action by $\op{Gal}(L/k)$, so $\mc F(k^{\mathrm{sep}})$ is a continuous module for $G_k\coloneqq\op{Gal}(k^{\mathrm{sep}}/k)$; thus, we have defined a functor $\op{Sh}(X_{\mathrm{\acute et}})\to\op{Mod}_\ZZ(G_k)$. There is also an inverse functor which sends $M\in\op{Mod}_\ZZ(G_k)$ to the sheaf $L\mapsto M^{\op{Gal}(k^{\mathrm{sep}}/L)}$, and one can check that these functors are equivalences. By comparing the notions of global sections on both sides, we see that
	\[\mathrm H^i_{\mathrm{\acute et}}(X;\mc F)=\mathrm H^i(G_k;\mc F(k^{\mathrm{sep}})).\]
\end{example}

\subsection{Brauer Groups of Schemes}
Throughout, we fix a scheme $X$.\footnote{One could even work with algebraic spaces or stacks, but we will not bother.} We would like to retell our story of the Brauer group of a field to work on general schemes. Accordingly, our first goal will be to redefine the notion of a central simple algebra.
\begin{definition}[Azumaya algebra]
	Fix a scheme $X$. Then an \textit{Azumaya algebra} $\mc A$ on $X$ is an $\mathcal O_X$-algebra $\mc A$ satisfying the following.
	\begin{listalph}
		\item The sheaf $\mc A$ is coherent over $\OO_X$.
		\item The algebra $\mc A$ is \'etale-locally isomorphic to a matrix algebra. Explicitly, there is an \'etale cover $\{U_i\}$ of $X$ and a tuple of ranks $\{r_i\}$ such that there are isomorphisms $\mc A|_{U_i}\cong M_{r_i}(\OO_{U_i})$ for each $i$.
	\end{listalph}
\end{definition}
\begin{remark}
	Here, (a) is analogous to requiring that a central simple algebra be finite-dimensional. The second condition is analogous to \Cref{prop:csa-grab-bag}(ii); we remark that (ii) implies that $\mc A$ is locally free.
\end{remark}
\begin{remark}
	If $X$ is connected, then it turns out that the rank of $\mc A$ is constant because it is locally free!
\end{remark}
\begin{remark} \label{rem:different-azumaya}
	The notion of an Azumaya algebra is older than the notion of the \'etale topology. Instead of the above, one can replace (b) that $\mc A$ is Zariski-locally free, nonzero at every point, and $\mc A\otimes\mc A\opp\cong\mc End(\mc A)$. This last condition comes from some non-commutative algebra.
\end{remark}
\begin{example}[quaternion algebra]
	Fix a scheme $X$ such that $2$ is a unit in $\OO_X$, and choose $f,g\in\OO_X(X)^\times$. Then
	\[\mc H_{f,g}\coloneqq\frac{\OO_X\langle i,j\rangle}{\left(i^2-f,j^2-g,ij+ji\right)}\]
	is an Azumaya algebra. This is analogous to \Cref{ex:quaternion-csa}.
\end{example}
We can now use these algebras to define a Brauer group, using a similar notion of equivalence.
\begin{definition}[Brauer group]
	Fix a scheme $X$. Then two Azumaya algebras $\mc A$ and $\mc A'$ are \textit{equivalent} if and only if there are locally free coherent $\OO_X$-modules $\mc E$ and $\mc E'$ of positive rank everywhere on $X$ such that
	\[\mc A\otimes_{\OO_X}\mc End_{\OO_X}(\mc E)\cong\mc A'\otimes_{\OO_X}\mc End_{\OO_X}(\mc E').\]
	Then we define the \textit{Azumaya Brauer group} $\op{Br}_{\op{Az}}(X)$ to be the collection of Azumaya algebras up to equivalence.
\end{definition}
\begin{remark}
	Similar to \Cref{rem:br-is-group}, the Brauer group is a group with operation given by the tensor product, and inverses given by the opposite Azumaya algebra.
\end{remark}
Alternatively, one could have tried to generalize the Brauer group by taking \'etale cohomology.
\begin{definition}[Brauer group]
	Fix a scheme $X$. Then the \textit{cohomological Brauer group} to be $\op{Br}X\coloneqq\mathrm H^2_{\mathrm{\acute et}}(X;\mathbb G_m)$.
\end{definition}
\begin{remark}
	A direct generalization of the argument from \Cref{prop:brauer-cohomology} provides a map $\op{Br}_{\op{Az}}X\to\op{Br}X$.
\end{remark}
One may hope that the natural map $\op{Br}_{\op{Az}}X\to\op{Br}X$ is an isomorphism, and this is the case in favorable circumstances.
\begin{theorem}[Gabber, de~Jong] \label{thm:gabber}
	Fix a scheme $X$. If $X$ is quasi-projective over an affine scheme, then the natural map $\op{Br}_{\op{Az}}X\to\op{Br}X$ is injective, and it surjects onto $(\op{Br}X)_{\mathrm{tors}}$.
\end{theorem}
\begin{remark}
	Gabber never published his proof of this result, so it can be difficult to find.
\end{remark}
\Cref{thm:gabber} tells us that we would like to know when $\op{Br}X$ is torsion.
\begin{proposition} \label{prop:function-field-injective}
	Fix a regular integral Noetherian scheme $X$. Then the map $\op{Br}X\to\op{Br}K(X)$ is injective.
\end{proposition}
\begin{proof}[Idea]
	One can think of elements of the Brauer group as sections of the Brauer--Severi variety. Under regularity assumptions, one can always do this splitting.
\end{proof}
\begin{remark}
	It follows from \Cref{prop:function-field-injective} that $\op{Br}X$ is torsion when $X$ is regular, integral, and Noetherian.
\end{remark}
\begin{remark}
	There are examples of singular varieties where $\op{Br}X$ fails to be torsion.
\end{remark}

\subsection{Some Computational Tools}
Here are some applications of \Cref{thm:gabber}.
\begin{example}
	Here is a silly example where $\op{Br}X$ has torsion: take $\bigsqcup_{i\in\NN}\Spec\QQ_2$. Then the cohomological Brauer group has an element $(1/2,1/3,1/4,\ldots)\in\prod_{i\in\NN}\QQ/\ZZ$, which is not torsion.
\end{example}
\begin{example}
	If $X$ is a nice curve over an algebraically closed field $k$, then $\op{Br}X$ is trivial. Indeed, $\op{Br}K(X)$ is trivial by Tsen's theorem, so this follows from \Cref{prop:function-field-injective}.
\end{example}
\begin{remark}
	In fact, if $X$ is a nice curve over a separably closed field $k$
\end{remark}
We can even go back and do some calculations for fields.
\begin{proposition} \label{prop:br-dvr}
	Fix a complete discrete valuation ring $R$ with residue field $k$. Then the natural map $\op{Br}R\to\op{Br}k$ is 
\end{proposition}
\begin{proof}
	We may pass to Azumaya algebras by \Cref{thm:gabber}. This is an application of Hensel's lemma. For injectivity, the main point is to check that if $\mc A_k$ is equivalent to a matrix algebra, then $\mc A$ is equivalent to a matrix algebra, which one does by finding idempotents in $k$ and using Hensel's lemma to lift them back to $\mc A$. For surjectivity, one uses some deformation theory. The idea is to successively lift the Azumaya algebra from $k$ to $R/\mf p$ and then $R/\mf p^2$ and so on, where here $\mf p$ is the maximal ideal of $R$. All these lifts turn out to be unique, so the lifting can be taken all the way to $R$.
\end{proof}
\begin{corollary}
	Fix a global field $k$. Then the image of the map
	\[\op{Br}k\to\prod_v\op{Br}k_v\]
	lands in $\bigoplus_v\op{Br}k_v$.
\end{corollary}
\begin{proof}
	For a finite set of places $S$ of $k$ including the infinite ones, let $\OO_{k,S}$ be the ring of $S$-integers, which is
	\[\OO_{k,S}=\{x\in k:v(x)\ge0\text{ for finite }v\in S\}.\]
	Now, there is a map $\OO_{k,S}\to\prod_{v\in S}k_{v}\times\prod_{v\notin S}\OO_v$, so $\op{Br}\OO_{k,S}$ embeds into $\prod_v\op{Br}k_v$ because $\op{Br}\OO_v=0$ by \Cref{prop:br-dvr}. By taking the direct limit, we see that
	\[\colim_S\op{Br}\OO_{k,S}\to\bigoplus_{v}\op{Br}k_v.\]
	To complete the proof, we should show that this map factors through $\op{Br}k$, which holds by passing to Azumaya algebras (via \Cref{thm:gabber}): indeed, any Azumaya algebra on $k$ is defined with finitely many elements (e.g., on a basis with some explicit multiplication law), which have finitely many denominators, so it can be lifted to $\op{Br}\OO_{k,S}$ for some potentially large $S$. It follows that the natural map $\colim_S\op{Br}\OO_{k,S}\to\op{Br}k$ is surjective, so we are done!
\end{proof}
In the sequel, it will be helpful to have a spectral sequence to compute our Brauer groups.
\begin{theorem}[Hochschild--Serre spectral sequence] \label{thm:h-s-spectral}
	Fix a nice $k$-variety $X$, and set $G\coloneqq\op{Gal}(k^{\mathrm{sep}}/k)$. Then there is a spectral sequence $E$ such that
	\[E_2^{pq}=\mathrm H^p\left(G;\mathrm H^q_{\mathrm{\acute et}}(X_{k^{\mathrm{sep}}};\mathbb G_m)\right)\Rightarrow\mathrm H^{p+q}_{\mathrm{\acute et}}(X;\mathbb G_m).\]
\end{theorem}
\begin{proof}
	This follows from Grothendieck's composition-of-functors spectral sequence 
\end{proof}
\begin{corollary}
	Fix a nice $k$-variety $X$, and set $G\coloneqq\op{Gal}(k^{\mathrm{sep}}/k)$. Let $\op{Br}_1X$ be the kernel of the map $\op{Br}X\to\op{Br}X_{k^{\mathrm{sep}}}$. Then there is an exact sequence
	\[0\to\op{Pic}X\to(\op{Pic}X_{k^{\mathrm{sep}}})^G\to\op{Br}k\to\op{Br}_1X\to\mathrm H^1(G;\op{Pic}X_{k^{\mathrm{sep}}})\to\mathrm H^3(G;k^{\mathrm{sep}\times}).\]
\end{corollary}
\begin{proof}
	This follows from the exact sequence associated to the spectral sequence, in addition to the following calculations.
	\begin{itemize}
		\item One has $\mathrm H^0_{\mathrm{\acute et}}(X_{k^{\mathrm{sep}}};\mathbb G_m)=k^{\mathrm{sep}\times}$ because the global sections of a projective variety are the constant functions.
		\item One has $\mathrm H^1_{\mathrm{\acute et}}(X_{k^{\mathrm{sep}}};\mathbb G_m)=\mathrm H^1_{\mathrm{Zar}}(X_{k^{\mathrm{sep}}};\mathbb G_m)=\op{Pic}X_{k^{\mathrm{sep}}}$. Indeed, these groups classify \'etale or Zariski $\mathbb G_m$-bundles, which one can show are all realized as line bundles, and $\op{Pic}X_{k^{\mathrm{sep}}}$ is the group of line bundles up to isomorphism.
		\item By definition, $\mathrm H^2_{\mathrm{\acute et}}(X_{k^{\mathrm{sep}}};\mathbb G_m)=\op{Br}X_{k^{\mathrm{sep}}}$ and similarly for $X$.
		\qedhere
	\end{itemize}
\end{proof}
\begin{remark}
	The cohomological proofs of class field theory show that $\mathrm H^3(G;k^{\mathrm{sep}\times})=0$ when $k$ is local or global.
\end{remark}
We close class with another criterion. For motivation, note that any normal integral Noetherian scheme $X$ admits an exact sequence
\[1\to\OO(X)^\times\to K(X)^\times\to\bigoplus_{x\in X^{(1)}}\ZZ,\]
where $X^{(1)}$ is the set of codimension-$1$ points. Namely, we are asserting that a meromorphic function $f\in K(X)^\times$ is a regular function if and only if it has vanishing valuation at all codimension-$1$ points. The analog of this for Brauer groups is as follows.
\begin{theorem} \label{thm:harthog-for-brauer}
	Fix a regular integral Noetherian scheme $X$. Then there is an exact sequence
	\[0\to\op{Br}X\to\op{Br}K(X)\to\bigoplus_{x\in X^{(1)}}\mathrm H^1(k(x);\QQ/\ZZ).\]
\end{theorem}
\begin{remark}
	We will not prove this or even define the maps on the right, but we will say that this recovers the left-exactness of \Cref{ex:cft-exact-seq} when $k$ is a function field.
\end{remark}
\begin{corollary}
	Fix a Zariski open over $\{U_i\}$ of a regular integral Noetherian scheme $X$. Then
	\[\op{Br}X=\bigcap_i\op{Br}U_i,\]
	where the intersection takes place in $\op{Br}K(X)$.
\end{corollary}
\begin{proof}
	To begin, use \Cref{thm:harthog-for-brauer} to present $\op{Br}X$ as elements of $\op{Br}K(X)$ which vanish in the various groups $\mathrm H^1(k(x);\QQ/\ZZ)$. Then the right-hand side is also presented as elements of $\op{Br}X$ which vanish in the various groups $\mathrm H^1(k(x);\QQ/\ZZ)$ where $x\in U_i$, so we recover the result by looping over all $U_i$s.
\end{proof}

\end{document}